\documentclass[twoside,11pt]{article}
\usepackage{jmlr2e}
\usepackage{booktabs}
\usepackage{tabularx}
\usepackage{array}
\usepackage{multirow}
\usepackage{path}

\jmlrheading{}{}{}{}{}{Herrera}

\ShortHeadings{XAI Metrics: Taxonomy and LIME--SHAP Comparison -- Supplementary}{Herrera}

\firstpageno{1}

\begin{document}

\appendix

\section*{Supplementary Material}
\label{sec:supplementary}

This supplementary document contains three tables referenced from the main
text: (S1) the exact evaluation rubric and prompts used in EXP4; (S2) the
LIME \texttt{kernel\_width} sensitivity probe results; and (S3) pairwise
feature correlations for the Adult Income dataset.

%% ────────────────────────────────────────────────────────────────────────────
\subsection*{Table S1 — EXP4 Evaluation Rubric and Judge Prompts}
\label{sec:supp_exp4_rubric}
%% ────────────────────────────────────────────────────────────────────────────

EXP4 employed three LLM judges each evaluating $n = 147$ paired SHAP/LIME
explanation outputs on seven semantic dimensions using a structured rubric.
The template was rendered per instance with the following fixed prompt
structure:

\medskip
\noindent\textbf{System instruction (delivered once per judge session):}
\begin{quote}\small
You are an expert AI evaluator specialising in Explainable AI (XAI) for
tabular data. Your task is to evaluate the quality of natural-language
explanation summaries produced for a binary classification model
(UCI Adult Income dataset, target: annual income $>$\$50K). Evaluate each
explanation on seven dimensions using a 1--5 integer Likert scale. Return
your response in strict JSON format with keys \texttt{scores} (object),
\texttt{rationales} (object), and optional \texttt{flags} (array). Do not
add commentary outside the JSON block.
\end{quote}

\medskip
\noindent\textbf{Per-instance user prompt (template variables in \{\{\}\}):}
\begin{quote}\small
\textbf{Dataset}: UCI Adult Income (Census Income, 14 features).\\
\textbf{Model}: \{\{model\_name\}\} (e.g., Random Forest, XGBoost).\\
\textbf{Explainer}: \{\{method\_name\}\} (SHAP or LIME).\\
\textbf{Predicted label}: \{\{prediction\_label\}\}. \textbf{True label}:
\{\{true\_label\}\}.\\[4pt]
\textbf{Explanation text}:\\
\{\{explanation\_text\}\}\\[4pt]
Rate the explanation on the seven dimensions below (1 = very poor,
5 = excellent). Return strict JSON.
\end{quote}

\begin{table}[h]
\centering
\small
\caption{EXP4 rubric: seven semantic evaluation dimensions, operational
definitions, and scoring anchors. Temperature = 0.0 for all judges.
Three primary condition: \texttt{hidden\_label} (standard), \texttt{label\_visible}
(bias probe), \texttt{rubric\_alt} (rubric sensitivity). Primary ICC
results come from the \texttt{hidden\_label} condition.}
\label{tab:s1_rubric}
\begin{tabularx}{\linewidth}{l p{4.2cm} p{5.6cm}}
\toprule
\textbf{Dimension} & \textbf{Operational definition} & \textbf{Scoring anchors} \\
\midrule
Clarity
  & Understandable organisation and wording for a non-expert reader.
  & 1 = incomprehensible; 3 = requires domain knowledge; 5 = clear to any reader \\[4pt]
Concision
  & Compact relative to useful information content; avoids unnecessary redundancy.
  & 1 = verbose or padded; 3 = adequate; 5 = maximally informative per word \\[4pt]
Actionability
  & Supports recourse, decision review, or explicit next steps for a user.
  & 1 = no actionable content; 3 = vague suggestions; 5 = specific, feasible actions \\[4pt]
Audit Usefulness
  & Supports inspection of model behaviour by an auditor or data scientist.
  & 1 = no audit value; 3 = partially informative about model logic; 5 = fully supports bias/error investigation \\[4pt]
Completeness
  & Provides sufficient local decision information to understand the specific prediction.
  & 1 = missing key drivers; 3 = covers main drivers; 5 = full local decision context \\[4pt]
Semantic Plausibility
  & Coherent with the instance context and general domain knowledge.
  & 1 = implausible or contradictory; 3 = plausible; 5 = highly consistent with domain expectations \\[4pt]
Overall Quality
  & Holistic semantic quality integrating all dimensions above.
  & 1 = very poor; 3 = acceptable; 5 = excellent \\
\bottomrule
\end{tabularx}
\end{table}

\noindent\textbf{Judge models and configuration:}
\begin{itemize}\itemsep0pt
  \item Judge A: \texttt{gpt-4o-mini} (OpenAI), temperature 0.0, max\_tokens 512.
  \item Judge B: \texttt{claude-3-haiku-20240307} (Anthropic), temperature 0.0, max\_tokens 512.
  \item Judge C: \texttt{gemini-1.5-flash} (Google), temperature 0.0, max\_tokens 512.
\end{itemize}

\noindent All judges used identical prompts; the rubric template is archived at
\path{src/prompts/templates/explanation_eval.j2} in the companion repository.

%% ────────────────────────────────────────────────────────────────────────────
\subsection*{Table S2 — LIME \texttt{kernel\_width} Sensitivity Probe}
\label{sec:supp_kw}
%% ────────────────────────────────────────────────────────────────────────────

Table~\ref{tab:s2_kw} reports LIME metrics under four \texttt{kernel\_width}
settings (RF model, seed 42, $N=100$, \texttt{num\_samples=1000},
\texttt{num\_features=10}). Fidelity and stability are computed only over
instances with at least one non-zero weight (``valid instances'').

\begin{table}[h]
\centering
\caption{LIME sensitivity to \texttt{kernel\_width} (RF, seed 42, $N=100$).
Fidelity = Pearson $r$ between $|w_i|$ and per-feature prediction drop.
Stability = mean pairwise cosine similarity under $T=15$ Gaussian
perturbations ($\sigma=0.1$). Reference condition: \texttt{kernel\_width=3.0}.}
\label{tab:s2_kw}
\begin{tabular}{lrcccc}
\toprule
\texttt{kernel\_width} & Valid instances & Fidelity & Stability & Sparsity & Cost (ms) \\
\midrule
0.25  & 0/100 (degenerate) & ---   & ---   & 0.000 & 30 \\
0.75  & 0/100 (degenerate) & ---   & ---   & 0.000 & 30 \\
3.0 \emph{(ref.)} & 96/100 & 0.518 & 0.000 & 0.087 & 30 \\
10.0  & 100/100            & 0.441 & 0.664 & 0.093 & 32 \\
\bottomrule
\end{tabular}
\end{table}

\noindent Kernel widths of 0.25 and 0.75 produce zero-weight explanations for
all 100 instances, because the exponential neighbourhood collapses in the
103-dimensional one-hot-encoded feature space. The reference value (3.0) is the
minimum viable setting. \texttt{kernel\_width=10.0} yields stable explanations
($0.664$) by averaging over a wider neighbourhood, at a cost of lower fidelity
($0.441$). These results confirm that LIME instability under the reference
configuration is not a calibration artifact: narrowing the kernel does not
improve stability; it degenerates the explanations. Widening the kernel
improves stability but at the cost of locality.

%% ────────────────────────────────────────────────────────────────────────────
\subsection*{Table S3 — Adult Income Feature Correlation Structure}
\label{sec:supp_corr}
%% ────────────────────────────────────────────────────────────────────────────

Table~\ref{tab:s3_corr} contextualises the out-of-distribution masking
discussion in the main text. Numeric--numeric correlations were computed as
Pearson $r$ on the raw (pre-preprocessing) training set; categorical--categorical
associations as Cramér's $V$ on the full dataset ($n = 48{,}842$).
Only pairs with $|r| > 0.10$ or $V > 0.15$ are shown.

\begin{table}[h]
\centering
\caption{Selected feature associations in the Adult Income dataset.
Pearson $r$ for numeric pairs; Cramér's $V$ for categorical pairs.
\texttt{education} and \texttt{education-num} are functionally identical
(one-to-one mapping; $V \equiv 1.0$). All values computed on the full
pre-split dataset ($n = 48{,}842$).}
\label{tab:s3_corr}
\begin{tabular}{llcc}
\toprule
Feature 1 & Feature 2 & Type & Association \\
\midrule
\multicolumn{4}{l}{\textit{Deterministic / near-perfect relationships}} \\
\texttt{education}     & \texttt{education-num}  & ordinal encode  & $V \equiv 1.0$ \\
\midrule
\multicolumn{4}{l}{\textit{Strong categorical associations ($V > 0.40$)}} \\
\texttt{relationship}  & \texttt{sex}            & Cramér's $V$    & 0.647 \\
\texttt{marital-status} & \texttt{relationship}  & Cramér's $V$    & 0.488 \\
\texttt{marital-status} & \texttt{sex}           & Cramér's $V$    & 0.460 \\
\texttt{occupation}    & \texttt{sex}            & Cramér's $V$    & 0.422 \\
\midrule
\multicolumn{4}{l}{\textit{Moderate categorical associations ($0.15 < V \leq 0.40$)}} \\
\texttt{workclass}     & \texttt{occupation}     & Cramér's $V$    & 0.210 \\
\texttt{education}     & \texttt{occupation}     & Cramér's $V$    & 0.190 \\
\texttt{occupation}    & \texttt{relationship}   & Cramér's $V$    & 0.171 \\
\midrule
\multicolumn{4}{l}{\textit{Numeric correlations ($|r| > 0.10$; all others $< 0.10$)}} \\
\texttt{education-num} & \texttt{hours-per-week} & Pearson $r$     & 0.144 \\
\texttt{capital-gain}  & \texttt{education-num}  & Pearson $r$     & 0.125 \\
\bottomrule
\end{tabular}
\end{table}

\noindent The four strongest associations (\texttt{relationship}--\texttt{sex},
\texttt{marital-status}--\texttt{relationship}, \texttt{marital-status}--\texttt{sex},
\texttt{occupation}--\texttt{sex}) all involve protected attributes, meaning that
independent feature masking during fidelity evaluation discards statistically
plausible joint values for these features. The functional duplicate
\texttt{education}/\texttt{education-num} means that any fidelity measurement
which masks one but not the other is measuring the effect of partial
information removal rather than true feature exclusion. These associations
do not invalidate the SHAP--LIME comparison---both methods are evaluated under
identical masking protocols---but they limit the interpretation of absolute
fidelity values as measures of causal feature importance.

\end{document}
